This computational example numerically fits cycloidal brachistochrone paths to
specified endpoints and checks the analytic travel-time result. It accompanies
CV17, \emph{The Brachistochrone Problem}.

\section{Endpoint solve}
For an endpoint $(X,Y)$, with $Y$ measured positively downward, write
\begin{equation}
 x=a(\theta-\sin\theta),\qquad
 y=a(1-\cos\theta).
\end{equation}
Eliminating $a$ gives the scalar nonlinear equation
\begin{equation}
 \frac{X}{Y}=\frac{\theta_B-\sin\theta_B}{1-\cos\theta_B}.
\end{equation}
The supplied Julia program solves this equation by bisection on
$0<\theta_B<2\pi$. In addition to the final root, it saves every bisection
bracket, midpoint, and residual so the numerical solve can be inspected rather
than treated as a black box. It then recovers
\begin{equation}
 a=\frac{Y}{1-\cos\theta_B},\qquad
 T=\theta_B\sqrt{\frac{a}{g}}.
\end{equation}

\section{CV17 reference case}
For $X=4\,\mathrm{m}$, $Y=1\,\mathrm{m}$ and
$g=9.81\,\mathrm{m\,s^{-2}}$, the numerical solve gives
\begin{equation}
 \theta_B\approx4.3760724130,\qquad
 a\approx0.7518727660\,\mathrm{m},\qquad
 T\approx1.2114965265\,\mathrm{s}.
\end{equation}
Because $\theta_B>\pi$, the path reaches a maximum depth
$2a\approx1.503745532\,\mathrm{m}$ before rising to the endpoint. The straight
segment requires about $1.86168\,\mathrm{s}$, so the cycloid is about $34.9\%$
faster for this example.

\begin{center}
\includegraphics{comparison.png}
\end{center}

\section{A second numerical calculation: travel-time convergence}
The exact cycloid time provides a useful target against which a direct numerical
quadrature can be tested. Divide the cycloid into $N$ short straight segments.
For segment $k$, let $\Delta s_k$ be its chord length and let $\bar y_k$ be the
midpoint depth. The numerical estimate is
\begin{equation}
 T_N=\sum_{k=1}^{N}\frac{\Delta s_k}{\sqrt{2g\bar y_k}}.
\end{equation}
The Julia resource saves this estimate for
$N=4,8,16,\ldots,16384$. The corresponding relative error decreases steadily as
the curve is refined. This is intentionally a separate calculation from the
closed-form identity $T=\theta_B\sqrt{a/g}$.

\section{Interactive explorer}
The static explorer coordinates four views of the same saved computation:
\begin{enumerate}
\item the cycloid overlaid on the straight chord, with a scrubber through saved
      position, time, and speed samples;
\item exact travel-time bars for the cycloid, straight line, and the idealized
      down-then-horizontal comparison;
\item the saved convergence history of $T_N$ toward the exact cycloid time; and
\item the saved bisection history for the endpoint root, including bracket
      width and endpoint-equation residual.
\end{enumerate}
The browser performs no endpoint solve and no travel-time integration. The
controls only select among Julia-generated samples already stored with the
resource.

\section{Reproducibility checks}
The Julia tests verify endpoint closure, the energy relation $v^2=2gy$, the
special $\theta_B=\pi$ case, scaling of geometrically similar problems, the
validity and contraction of every bisection bracket, and monotone convergence
of the independent segment-time calculation.

The \PMlinkexternal{interactive explorer and downloadable Julia project}
{https://images.physicslibrary.org/examples/brachistochrone-cycloid/index.html}
include source, tests, trajectory CSVs, root history, convergence data, plots,
and provenance.
